\begin{proofbox}
  \textbf{\textcolor{CProof}{方法一：先用一个例子看懂（直观验证）}}

  把 $y=x^2$ 左移 $2$ 格。原顶点 $(0,0)$ 变成 $(-2,0)$，原点 $P(1,1)$ 变成 $Q(-1,1)$。
  看新式 $y=(x+2)^2$：
  \[
     x=-2\Rightarrow y=0,\qquad x=-1\Rightarrow y=1.
  \]
  这两个点帮助我们发现规律，但仅验证两个点并不能证明整条图像都平移了。
  再任取原抛物线上的点 $(t,t^2)$，左移 $2$ 后是 $(t-2,t^2)$。
  代入新式得到 $(t-2+2)^2=t^2$，因此每个原图像上的点都能对应到新图像；反过来，新图像任一点 $(x,(x+2)^2)$ 向右移 $2$，也回到原抛物线上。
  这就验证了抛物线的平移；推广到任意函数见方法二。
  \begin{center}
  \begin{tikzpicture}[x=0.9cm,y=0.75cm,>=stealth,font=\small]
    \draw[->,gray] (-3.7,0)--(1.95,0) node[right] {$x$};
    \draw[->,gray] (0,-0.2)--(0,2.5) node[above] {$y$};
    \draw[blue!75!black,thick,domain=-1.45:1.45,samples=40] plot (\x,{\x*\x});
    \draw[teal!75!black,thick,domain=-3.45:-0.55,samples=40] plot (\x,{(\x+2)*(\x+2)});
    \fill[blue!75!black] (0,0) circle (2pt);
    \fill[teal!75!black] (-2,0) circle (2pt);
    \draw[->,teal!75!black,thick] (-0.2,-0.58)--(-1.8,-0.58)
      node[midway,below] {$2$ 格};
    \node[below] at (-2,0) {$-2$};
    \node[below right] at (0,0) {$0$};
  \end{tikzpicture}

  \end{center}

  \medskip
  \textbf{\textcolor{CProof}{方法二：追踪任意一点（严格推导）}}

  设 $a>0$，原图像上任取一点 $P(u,f(u))$。
  \begin{description}[style=nextline,leftmargin=2.6em,labelsep=0.5em,itemsep=0.4em]
    \item[\textbf{左移：}] $P$ 左移 $a$ 后变成 $Q(u-a,f(u))$。若把新点的横坐标记为 $x=u-a$，则 $u=x+a$。
      所以新点的高度 $y=f(u)=f(x+a)$，即
      \[\boxed{g(x)=f(x+a)}.\]
    \item[\textbf{右移：}] $P$ 右移 $a$ 后变成 $R(u+a,f(u))$。若新点横坐标为 $x=u+a$，则 $u=x-a$。
      所以 $y=f(u)=f(x-a)$，即
      \[\boxed{g(x)=f(x-a)}.\]
  \end{description}
  因为 $P$ 是\textbf{任意一点}，这个规律不只对抛物线成立，对任意函数图像都成立。
  \begin{center}
  \begin{tikzpicture}[x=1cm,y=.85cm,>=stealth,font=\small]
    \draw[->,gray] (-1,-.6)--(10.2,-.6) node[right] {$x$};
    \draw[densely dashed,gray!70] (.5,1.1)--(8.5,1.1);
    \fill[teal!75!black] (.5,1.1) circle(2.4pt);
    \fill[blue!75!black] (4.5,1.1) circle(2.4pt);
    \fill[orange!85!black] (8.5,1.1) circle(2.4pt);
    \node[above,align=center,teal!75!black] at (.5,1.25) {$Q$\\$(u-a,f(u))$};
    \node[above,align=center,blue!75!black] at (4.5,1.25) {$P$\\$(u,f(u))$};
    \node[above,align=center,orange!85!black] at (8.5,1.25) {$R$\\$(u+a,f(u))$};
    \draw[->,teal!75!black,thick] (4.2,.5)--(.8,.5) node[midway,below] {左移 $a$};
    \draw[->,orange!85!black,thick] (4.8,.5)--(8.2,.5) node[midway,below] {右移 $a$};
    \node at (4.5,2.85) {同一高度};
\end{tikzpicture}

  \end{center}

  \medskip
  \textbf{\textcolor{CProof}{方法三：反过来验证公式确实表示平移}}

  先不看图，直接设 $g(x)=f(x+a)$。若原函数在 $u$ 处的高度为 $f(u)$，
  新函数在 $x=u-a$ 处的高度就是
  \[
     g(u-a)=f((u-a)+a)=f(u).
  \]
  也就是说，\textbf{每个点都向左走了 $a$ 格，高度完全不变}，这正是图像左移。
  同样，若 $g(x)=f(x-a)$，则
  \[
    g(u+a)=f((u+a)-a)=f(u),
  \]
  说明它是右移。反过来，若新图像上的点是 $(x,g(x))$，左移公式对应的原点为 $(x+a,f(x+a))$，右移公式对应的原点为 $(x-a,f(x-a))$。
  因而在相应的平移后定义域上，新旧点一一对应，不会多点或漏点。
  \begin{center}
  \begin{tikzpicture}[x=1cm,y=.85cm,>=stealth,font=\small]
    \draw[->,gray] (-1,-.6)--(10.2,-.6) node[right] {$x$};
    \draw[densely dashed,gray!70] (.5,1.1)--(8.5,1.1);
    \fill[teal!75!black] (.5,1.1) circle(2.4pt);
    \fill[blue!75!black] (4.5,1.1) circle(2.4pt);
    \fill[orange!85!black] (8.5,1.1) circle(2.4pt);
    \node[above,align=center,teal!75!black] at (.5,1.25) {$Q$\\$(u-a,f(u))$};
    \node[above,align=center,blue!75!black] at (4.5,1.25) {$P$\\$(u,f(u))$};
    \node[above,align=center,orange!85!black] at (8.5,1.25) {$R$\\$(u+a,f(u))$};
    \draw[->,teal!75!black,thick] (4.2,.5)--(.8,.5) node[midway,below] {左移 $a$};
    \draw[->,orange!85!black,thick] (4.8,.5)--(8.2,.5) node[midway,below] {右移 $a$};
    \node at (4.5,2.85) {同一高度};
\end{tikzpicture}

  \end{center}
  \textbf{反向读图：}从 $Q$ 向右走 $a$ 回到 $P$，从 $R$ 向左走 $a$ 也回到 $P$；每一步高度都不变。

  \textbf{\textcolor{CProof}{结论回扣：}}点的横坐标向左减 $a$，而写新函数必须找回原横坐标，所以\textbf{括号中的 $x$ 要加 $a$}。
\end{proofbox}
